\documentclass[11pt]{article}
\usepackage[T1]{fontenc}
\usepackage{lmodern}
\usepackage{amsmath,amssymb}
\usepackage[colorlinks=true,linkcolor=blue,urlcolor=blue]{hyperref}
\usepackage[margin=1in]{geometry}
\title{Huybrechts: Compact Hyperk\"ahler Manifolds\\Reading Seminar 1}
\date{}
\begin{document}
\maketitle

\noindent
Daniel Huybrechts, \emph{Compact Hyperk\"ahler Manifolds:
Basic Results}, \href{https://arxiv.org/abs/alg-geom/9705025}
{arXiv:alg-geom/9705025} (1997).
Reading: Introduction and Sections~1, 5--7.

{\small\tableofcontents}

\section{Historical and bibliographical background}

\subsection{Two origins of hyperk\"ahler geometry}
In Riemannian geometry, a hyperk\"ahler manifold
of real dimension $4n$ has holonomy contained in
$\operatorname{Sp}(n)$. Equivalently, its tangent
bundle carries three parallel complex structures
$I,J,K$ satisfying the quaternion relations.
This viewpoint grows out of Hamilton's quaternions
(1843) and Berger's holonomy classification (1955).
A second origin lies in the study of K3 surfaces:
compact complex surfaces with a holomorphic
symplectic form, whose Hodge structures and
period maps make their complex geometry unusually
accessible. The modern subject joins these two
traditions and generalizes K3 surfaces to higher
complex dimensions.

\subsection{1955--1978: Ricci-flat metrics and complex geometry}
Calabi conjectured the existence of a unique
Ricci-flat K\"ahler metric in each K\"ahler class
when the first Chern class vanishes
(\emph{On K\"ahler manifolds with vanishing
canonical class}, 1955); Yau proved this
conjecture (\emph{Comm. Pure Appl. Math.}
\textbf{31} (1978), 339--411).
Consequently, a simply connected compact
K\"ahler manifold with a unique (up to scale)
nondegenerate holomorphic two-form admits
a metric of holonomy $\operatorname{Sp}(n)$.
The name \emph{irreducible holomorphic symplectic}
(IHS) emphasizes the complex-geometric description;
\emph{irreducible hyperk\"ahler} emphasizes the metric
description. See Huybrechts \S\S1.1--1.3.

\subsection{K3 surfaces, deformation theory and Torelli}
Kodaira--Spencer's deformation theory (1958--1960)
made it possible to vary complex structures
systematically. The K3 global Torelli theorem
showed, roughly, that a K3 surface can be recovered
from the Hodge structure on $H^2$ together with the
intersection form, subject to the relevant marking
and K\"ahler-cone conditions; see Burns--Rapoport,
\emph{On the Torelli Problem for K\"ahlerian K3
Surfaces} (1975), Huybrechts [12].
The associated \emph{period} records the line
$H^{2,0}(X)\subset H^2(X,\mathbb C)$ under a
chosen cohomological marking.
Bogomolov's work (1974, 1978, [6,7])
placed holomorphic symplectic manifolds among
compact K\"ahler manifolds with $c_1=0$;
he also established unobstructedness of their
deformations.
These ideas suggested studying periods and moduli
for higher-dimensional analogues of K3 surfaces.

\subsection{1983: Beauville establishes the modern theory}
Beauville, \emph{Vari\'et\'es k\"ahleriennes dont la
premi\`ere classe de Chern est nulle},
\emph{J. Differential Geom.} \textbf{18} (1983),
755--782 [2], exhibited two fundamental series
of IHS manifolds: Hilbert schemes $S^{[n]}$
of points on K3 surfaces and generalized Kummer
varieties associated with abelian surfaces.
He also introduced a quadratic form $q_X$
on $H^2(X,\mathbb Z)$, generalizing the
intersection pairing of a K3 surface, and
proved \emph{local Torelli}: nearby marked
complex structures are parametrized locally
by their periods. This is the starting point
of Huybrechts \S\S1.9, 1.15 and 2.

\subsection{1984--1996: examples and cohomological structure}
Mukai constructed holomorphic symplectic forms
on moduli spaces of sheaves over K3 and abelian
surfaces (\emph{Invent. Math.} \textbf{77} (1984),
101--116 [39]). Beauville--Donagi (1985, [4])
studied varieties of lines on cubic fourfolds;
Hitchin--Karlhede--Lindstr\"om--Ro\v cek (1987,
[27]) developed the metric and hyperk\"ahler
quotient approach. Fujiki proved the
intersection formula, for $\dim_{\mathbb C}X=2n$,
\[
\int_X\alpha^{2n}=c_Xq_X(\alpha)^n
\]
(\emph{Adv. Stud. Pure Math.} \textbf{10}
(1987), 105--165 [21]). Together with
Tian's (1987, [48]) and Todorov's
(1989, [49]) extension of unobstructedness
to compact K\"ahler manifolds with trivial
canonical bundle---now commonly called the
Bogomolov--Tian--Todorov (BTT) theorem--- and the work of
Looijenga--Lunts and Verbitsky on cohomology
([36,51]), this explained why $H^2$ could
control geometry in higher dimensions.
But Debarre's counterexample (1984, [15])
showed that identical periods need not
imply isomorphic IHS manifolds.

\subsection{1997--1999: Huybrechts' basic results}
Huybrechts' present paper (1997 preprint;
\emph{Invent. Math.} \textbf{135} (1999))
addresses the global questions left by
Beauville's local theory: when is an IHS
manifold projective (\S3), how are birational
models related (\S4), what determines the
K\"ahler cone (\S\S5--7), and does every
admissible period arise (\S8)?
He proves global \emph{surjectivity} of
the period map, without claiming that
each period determines a unique manifold.
His \S10 identifies the remaining Torelli
and classification problems.
The projectivity theorem requires the
2003 erratum (\emph{Invent. Math.}
\textbf{152}, 209--212).

\subsection{1999--2005: more deformation types and cones}
O'Grady constructed further IHS deformation
types in complex dimensions ten and six
(\emph{J. Reine Angew. Math.} \textbf{512}
(1999), 49--117; \emph{J. Algebraic Geom.}
\textbf{12} (2003), 435--505).
Huybrechts refined the description of
K\"ahler cones (\emph{Math. Ann.}
\textbf{326} (2003), 499--513), and
Demailly--P\u{a}un obtained a general
numerical characterization of K\"ahler
cones (\emph{Ann. of Math.} \textbf{159}
(2004), 1247--1274).

\subsection{2013: Verbitsky's global Torelli theorem}
For marked IHS manifolds, the \emph{period map}
assigns the line $H^{2,0}$ in a fixed
lattice. Two points of a non-Hausdorff
moduli space are called \emph{inseparable}
if they cannot be distinguished by disjoint
open neighborhoods. Huybrechts had shown
that inseparable marked points give
birational manifolds (Theorem~4.3), and
that birational \emph{projective} IHS
manifolds yield inseparable points for
suitable markings (Theorem~4.6').
More precisely, with suitable markings,
\[
\begin{gathered}
X\dashrightarrow X'\text{ birational},\quad X,X'\text{ projective}
\\[2pt]
\Longrightarrow\ (X,\varphi),(X',\varphi')\text{ inseparable}
\\[2pt]
\Longrightarrow\ \operatorname{Per}(X,\varphi)=\operatorname{Per}(X',\varphi').
\end{gathered}
\]
The missing implication was
\emph{same period $\Rightarrow$ inseparability};
extending Theorem~4.6 to nonprojective
manifolds would not establish it.
Verbitsky, \emph{Mapping class group and a global
Torelli theorem for hyperk\"ahler manifolds},
\emph{Duke Math. J.} \textbf{162} (2013),
2929--2986, proved that the period map
becomes an isomorphism on each appropriate
\emph{birational Teichm\"uller component},
obtained by identifying inseparable
complex structures on a fixed smooth
manifold modulo isotopy.
For $I,I'$ in the same connected component
of HK Teichm\"uller space,
\[
\begin{gathered}
\operatorname{Per}(I)=\operatorname{Per}(I')
\ \Longrightarrow\ I,I'\text{ inseparable}
\\[2pt]
\Longrightarrow\ (M,I)\dashrightarrow(M,I')\text{ birational}.
\end{gathered}
\]
Connected-component
and monodromy qualifications are essential;
see D.~Huybrechts,
\emph{A global Torelli theorem for hyperk\"ahler
manifolds (after M.~Verbitsky)},
\href{https://arxiv.org/abs/1106.5573}
{arXiv:1106.5573} (2011).

\subsection{Recent developments: monodromy and walls}
The distinction between a Hodge isometry
and one induced by a geometric transformation
makes monodromy central to global Torelli.
Markman's
\href{https://arxiv.org/abs/1101.4606}
{survey} organizes these results.
Amerik--Verbitsky, \emph{Rational curves
on hyperk\"ahler manifolds} (\emph{IMRN},
2015), developed the theory of MBM
classes, whose orthogonal hyperplanes
describe walls separating K\"ahler
and birational chambers.

\subsection{Current open problems (October 2026)}
The classification of IHS deformation
types remains open: for instance, is
their number finite in each dimension,
and what other types exist beyond
$K3^{[n]}$, generalized Kummer, OG6
and OG10? Further problems involve
explicit birational models and monodromy,
the existence and structure of Lagrangian
fibrations, and degenerations and
compactifications of moduli.
Studying limits
of generalized Kummer varieties arising
from degenerating abelian surfaces is
one concrete example of the broader
degeneration problem.

\section{Contributions of Huybrechts' 1997 paper}

\noindent
\textbf{This is an original research paper, not merely a survey.}
Its main results include:
\begin{itemize}
\item \textbf{Theorem~3.11:} $X$ is projective
if and only if some line bundle $L$ satisfies
$q_X(c_1(L))>0$. The correct proof is in
the 2003 erratum, not the original proof.
\item \textbf{Theorem~4.6 and Corollary~4.7:}
birational \emph{projective} irreducible symplectic
manifolds deform in families isomorphic away from
the central fibre; hence they are diffeomorphic.
\item \textbf{Proposition~5.1:} an analogue of
the transitivity of the K3 Weyl action, with
a K\"ahler class realized on a suitable
non-separated birational model.
\item \textbf{Corollary~6.5 and Theorem~7.1:}
ample/K\"ahler cone criteria, weaker than
the explicit K3 $(-2)$-curve criterion.
\item \textbf{Theorem~8.1:} the global period map
$\mathcal P:\mathcal M_\Gamma^0\to\mathcal D_\Gamma$
is \emph{surjective} for every nonempty connected
component of the marked moduli space.
This says every admissible period occurs,
not that it occurs uniquely.
\end{itemize}

\medskip\noindent
\textbf{Guide to the paper.}
Section~1 supplies the preliminaries: hyperk\"ahler
geometry, the Beauville--Bogomolov form,
deformations, and periods; Section~2 develops
the examples. The main new results begin with
the projectivity criterion in Section~3 and
the birational deformation theorem in Section~4.
Sections~5--7 study how K\"ahler and ample
cones distinguish birational models. These
results lead to the central global statement,
Theorem~8.1 in Section~8: the period map is
surjective. Section~9 explores consequences,
and Section~10 discusses the Torelli and
classification questions left open.

\section{Preliminaries (Section 1): a reading guide}

\noindent
\textbf{Big picture.}
Section~1 assembles the language needed for the rest of
the paper: a holomorphic symplectic manifold has a
hyperk\"ahler metric, a quadratic form on $H^2$, a
smooth deformation space, and a period map.
The central theme is that degree-two cohomology
governs both deformations and K\"ahler geometry.

\begin{itemize}
\item Definitions and hyperk\"ahler rotation (1.1--1.5).
\item Cohomology and the Beauville--Bogomolov form (1.6--1.11).
\item Kodaira--Spencer and unobstructed deformations (slightly off-topic).
\item Deforming a line bundle (1.14) [understanding in-progress].
\item Markings and the period map (1.15--1.18) [understanding in-progress].
\item Twistor families and Hodge loci (1.13, 1.16--1.17) [understanding in-progress].
\end{itemize}

\par\smallskip\noindent
\textbf{Definitions and hyperk\"ahler rotation
(1.1--1.5).}
An \emph{irreducible symplectic manifold} is a
simply connected compact K\"ahler manifold with
$H^0(X,\Omega_X^2)=\mathbb C\sigma$, where $\sigma$
is everywhere nondegenerate. Yau's theorem relates
this to irreducible hyperk\"ahler metrics.
The parallel complex structures
$aI+bJ+cK$ ($a^2+b^2+c^2=1$) form a
sphere; this is the origin of twistor families.
The holomorphic symplectic form identifies
$T_X\simeq\Omega_X^1$.

\par\smallskip\noindent
\textbf{Cohomology and the Beauville--Bogomolov
form (1.6--1.11).}
Fujiki's Lefschetz-type decompositions,
the holomorphic symplectic form, and
Serre duality yield cohomological symmetries
and perfect pairings. The key object is the
primitive integral quadratic form
$q_X$ on $H^2(X,\mathbb Z)$, of signature
$(3,b_2-3)$, satisfying Fujiki's formula
$\int_X\alpha^{2n}=c_Xq_X(\alpha)^n$.
The \emph{positive cone} $\mathcal C_X$ is
the component of
$\{\alpha\in H^{1,1}(X,\mathbb R):
q_X(\alpha)>0\}$ containing the K\"ahler
cone $\mathcal K_X$; these cones need
not coincide.

\par\smallskip\noindent
\textbf{Kodaira--Spencer and unobstructed deformations (slightly off-topic).}
For a deformation $\pi:\mathcal X\to B$ with
$X=\pi^{-1}(0)$, restriction of the relative
tangent sequence gives
\[
0\to T_X\to T_{\mathcal X}|_X
\to T_0B\otimes\mathcal O_X\to0.
\]
The connecting homomorphism is the
Kodaira--Spencer map
$\rho:T_0B\to H^1(X,T_X)$; for the
semiuniversal family it is an isomorphism.
Contraction with the holomorphic symplectic
form identifies $T_X\simeq\Omega_X^1$,
and Dolbeault gives
$H^1(X,T_X)\simeq H^{1,1}(X)$.
Bogomolov--Tian--Todorov implies that
the deformation space is smooth: obstruction
classes vanish even though the obstruction
group $H^2(X,T_X)$ need not vanish.

\par\smallskip\noindent
\textbf{Deforming a line bundle (1.14).}
A deformation of the pair $(X,L)$ consists
of a family $\pi:\mathcal X\to B$ and a
holomorphic line bundle $\mathcal L$ on
the total space satisfying
$\mathcal L|_X\simeq L$.
Smooth local triviality transports
the topological class $c_1(L)$ to nearby
fibres, but it need not remain the first
Chern class of a \emph{holomorphic} line
bundle. By the Lefschetz $(1,1)$-theorem,
it must stay of Hodge type $(1,1)$
(because the Lefschetz $(1,1)$-theorem says
every integral $(1,1)$-class is the class
of a line bundle).
The Beauville--Bogomolov form is compatible
with the Hodge decomposition: a real class
$\alpha\in H^2(X_t,\mathbb R)$ is of type
$(1,1)$ precisely when it is orthogonal
to $H^{2,0}(X_t)=\mathbb C\sigma_t$.
Thus the Lefschetz $(1,1)$ condition on
the transported class $c_1(L)$ becomes
\[
q_X(c_1(L),\sigma_t)=0.
\]
Here cohomology is identified locally across
the family; $c_1(L)$ stays fixed while
$\sigma_t$ varies. Since local Torelli makes
the period $[\sigma_t]$ a local coordinate
on the deformation space, this is one
holomorphic equation in period coordinates.
For $c_1(L)\ne0$ it cuts out a smooth
hypersurface of the local deformation space,
of dimension $h^{1,1}(X)-1$.
Locally the fibrewise holomorphic bundles
assemble into an extension $\mathcal L$.

\par\smallskip\noindent
\textbf{Markings and the period map (1.15--1.18).}
Cohomology groups of different fibres are
not literally the same vector space.
A \emph{marking} is an isometry
\[
\varphi:(H^2(X,\mathbb Z),q_X)
\xrightarrow{\ \sim\ }(\Gamma,q_\Gamma)
\]
with a fixed lattice. It lets us compare
all periods in the same projective space:
\[
\mathcal P(X,\varphi)
=[\varphi(\sigma_X)]\in
\mathcal Q_\Gamma
=\{[x]\in\mathbb P(\Gamma_\mathbb C):
q_\Gamma(x)=0,\ q_\Gamma(x,\bar x)>0\}.
\]
Beauville's local Torelli theorem says
the period map is a local isomorphism.
The moduli space of marked manifolds
can be non-Hausdorff; it is the domain
of the global period map in
Huybrechts' Theorem~8.1.

\par\smallskip\noindent
\textbf{Twistor families and Hodge loci (1.13, 1.16--1.17).}
A hyperk\"ahler metric gives the sphere
of induced complex structures, producing
a twistor family over $\mathbb P^1$.
The paper also examines the loci where
a fixed cohomology class remains of type
$(1,1)$ and the behaviour of curves and
divisors along twistor families; these
constructions prepare the later study
of K\"ahler cones and birational models.

\end{document}
